\section{Propósito, contexto y alcance}

\subsection{Contexto del desarrollo}

El presente trabajo desarrolla una representación computacional de un motor síncrono de imanes permanentes (PMSM) Siemens 1FT7087 con el propósito de establecer una plataforma reproducible para el desarrollo y evaluación de estrategias de control.

El desarrollo comprende la caracterización de la planta, la investigación y análisis de la información disponible del motor, la formulación de relaciones físicas y matemáticas, la interpretación, conversión y derivación de parámetros, la parametrización del modelo, su implementación computacional y la realización de experimentos reproducibles.

La arquitectura fue desarrollada y validada inicialmente mediante Simulink, constituyendo una referencia para la posterior implementación reproducible del Digital Twin. La implementación en Python permite establecer un entorno independiente para la ejecución, registro y repetición de los experimentos bajo configuraciones explícitas.

\subsection{Problemática de ingeniería}

El desarrollo de estrategias avanzadas de control requiere una representación de la planta que mantenga una relación coherente entre sus parámetros, variables de estado, entradas, salidas, condiciones de operación y resultados experimentales.

En este contexto, la problemática no se limita a disponer de los datos nominales del motor, sino a establecer una representación computacional trazable en la que la información de la planta pueda relacionarse con su formulación física y matemática, su implementación y las condiciones bajo las cuales se obtienen los resultados.

Esta representación debe permitir, además, utilizar una misma planta de referencia para desarrollar y evaluar diferentes estrategias de control, evitando que las diferencias observadas entre estrategias se atribuyan a cambios no controlados en la representación de la planta o en las condiciones experimentales.

\subsection{Propósito del Digital Twin}

El propósito del Digital Twin es proporcionar una representación computacional específica del motor Siemens 1FT7087 que sirva como plataforma común para el desarrollo y evaluación de estrategias de control.

La planta PMSM implementada como M4\_BASE constituye el núcleo funcional actual de esta representación. Su formulación mediante parámetros concentrados proporciona el nivel de representación requerido para el objetivo actual de desarrollo y evaluación del control, manteniendo explícitas las relaciones entre los parámetros adoptados, las variables del modelo y las condiciones de operación.

El nivel de fidelidad del Digital Twin no se establece por la incorporación de la mayor cantidad posible de fenómenos o parámetros, sino de acuerdo con la información requerida por el objetivo de ingeniería. La representación puede ampliarse de forma trazable cuando una estrategia de control o un fenómeno de interés requiera información adicional.

\subsection{Alcance del reporte}

El reporte comprende la caracterización y trazabilidad de la planta PMSM, la definición de la arquitectura del Digital Twin, la formulación e implementación del modelo físico M4\_BASE, la definición de la interfaz entre planta y control, y el establecimiento de un controlador FOC como referencia experimental.

También se documentan la metodología experimental y de reproducibilidad, los ensayos realizados y sus resultados, los criterios de validación utilizados y la distinción entre consistencia computacional, evidencia funcional, reproducibilidad y validación física.

Finalmente, se establece la ruta de evolución hacia estrategias de control avanzado y la forma en que puede ampliarse la fidelidad de la representación de la planta de acuerdo con las necesidades futuras del desarrollo.

\subsection{Metodología general de desarrollo}

El desarrollo sigue una cadena de trabajo orientada a mantener la trazabilidad entre la planta real, su representación computacional y la evidencia experimental obtenida:

\begin{center}
\textit{investigación}
$\rightarrow$
\textit{análisis}
$\rightarrow$
\textit{cálculo y derivación}
$\rightarrow$
\textit{parametrización}
$\rightarrow$
\textit{formulación}
$\rightarrow$
\textit{implementación}
$\rightarrow$
\textit{experimentación}
$\rightarrow$
\textit{evaluación}.
\end{center}

Esta secuencia permite diferenciar los datos de origen de los parámetros derivados, las formulaciones adoptadas, las decisiones de implementación y los resultados obtenidos mediante experimentación computacional.

La metodología se mantiene como referencia común durante la evolución del Digital Twin, de manera que las modificaciones de la representación de la planta o de la estrategia de control puedan conservar condiciones experimentales identificables y resultados trazables.
